% !TeX root = main.tex
{\fontsize{12}{14}\selectfont\tenantselection}
\subsection*{Normal lighting and local controls}
The lighting design uses 92 PIL PLATO 600 MultiPower 3117955 panels on the 600 mm ceiling grid: sixty 22 W, twenty-eight 28 W and four 38 W, totalling 2256 W. Separate circuits serve the enclosed rooms and distributed main-office zones, with a maintained mean illuminance target of 500 lx adopted for the design. The connected lighting density is 7.301 W/m$^2$ using the adopted 309 m$^2$ office area.

Seven passive manual switches S-A--S-G provide local control of their corresponding circuits. These switches must have a minimum rating of 20 AX, an LED switching capacity of at least 600 W and make/break/inrush compatibility with the selected drivers. Since their phases differ, S-A/B/C use separate enclosures. No borrowed neutral is allowed. Each route passes through its switch and includes a 2 m DB rise plus 1.5 m down and 1.5 m back at the switch, giving an assumed vertical allowance of 5 m per lighting radial. Passive controls add no operating power.
\begin{center}\begin{tabular}{lrrrrl}\toprule
Circuit & Lamps & W & Phase & RCBO & Control\\\midrule

LC-A & 18 & 402 & L1 & B10 & S-A \\
LC-B & 17 & 392 & L2 & B10 & S-B \\
LC-C & 18 & 396 & L3 & B10 & S-C \\
LC-D & 11 & 242 & L3 & B10 & S-D \\
LC-E & 14 & 392 & L2 & B10 & S-E \\
LC-F & 10 & 280 & L1 & B10 & S-F \\
LC-G & 4 & 152 & L3 & B10 & S-G \\

\bottomrule\end{tabular}\end{center}
On the plan, the A/B/C numeric labels correspond to the LED IDs and circuits LC-A/B/C. D01--D11 correspond to PN01--PN11, E82--E96 excluding E90 to the northeast LEDs, F01--F10 to UE01--UE10, and G01--G04 to SE01--SE04.
% !TeX root = main.tex
\subsection*{Lighting calculation}
To assess lighting performance, the calculation uses manufacturer LDT intensities given in cd/klm, with each profile scaled by its selected flux: 2843.45 lm at 22 W, 3555.30 lm at 28 W and 4685.55 lm at 38 W. Illuminance is then calculated at each workplane grid cell, with spacing at most 0.2 m, as
\[
E=0.80\sum_j\frac{I_j(C,\gamma)\cos\gamma}{d_j^2}.
\]
A luminaire-to-workplane height of 1.9 m is adopted, corresponding to a 2.7 m ceiling and a 0.8 m workplane. The maintenance factor of 0.80 assumes periodic cleaning and replacement. Room boundaries use the adopted plan envelopes; each room receives light from its own fittings. Only direct light is included in this calculation.
\begin{center}\begin{tabular}{lrrrr}\toprule
Zone & Lamps & Area m$^2$ & Mean lx & Min/mean\\\midrule
Open office & 67 & 231.87 & 592.9 & 0.136\\
Pantry & 11 & 29.54 & 546.3 & 0.390\\
Upper east room & 10 & 33.68 & 588.5 & 0.267\\
South east room & 4 & 13.07 & 662.6 & 0.175\\
\bottomrule\end{tabular}\end{center}
On this basis, all four mean illuminances exceed the adopted 500 lx target. To provide mounting-height tolerance, the four south-east panels use the 38 W setting. At 2.2 m above the workplane, the respective means are 573.0, 505.8, 545.0 and 589.2 lx; the south-east room retains a 17.8\% margin.

Min/mean ratios include perimeter cells. These direct-light results do not establish task-area uniformity or glare performance.

\subsection*{Small power and fixed equipment}
Small-power provision includes nine four-socket floor boxes, three twin wall sockets and the fixed connection FS01. Each floor-box/socket assembly is individually 13 A fused. The demand calculation uses a 13 A continuous capacity per assembly; FS01 supplies WH01, a STIEBEL ELTRON SHC10 GB Eltron 235232 undersink water heater. Its 1.4 kW at 220 V gives 6.364 A at unity power factor; the selected protection is B10 with a 10 A double-pole FCU. Building central AC is supplied as central plant and not duplicated in this office DB.

P1--P4 are A1 rings with fused floor-box connections. P5/P6/P8 are A2 radials, while P7 is a fixed radial. Both P6 and P8 serve the PANTRY identified on the source plan, with FB09 and SO01 connected to P6. P8 is a separate socket group for convenient isolation. Kitchen sockets are separated from ordinary office circuits under EMSD6E(3)(a).
\begin{center}\begin{tabular}{llllr}\toprule
Circuit & Points & Phase & Topology & Protection\\\midrule

P1 & FB01, FB02 & L1 & A1 ring & B32 \\
P2 & FB03, FB04 & L1 & A1 ring & B32 \\
P3 & FB05, FB06 & L2 & A1 ring & B32 \\
P4 & FB07, FB08 & L2 & A1 ring & B32 \\
P5 & SO02 & L1 & A2 radial & B32 \\
P6 & FB09, SO01 & L3 & A2 radial & B32 \\
P7 & FS01/WH01 & L3 & Fixed radial & B10 \\
P8 & SO03 & L3 & A2 radial & B32 \\

\bottomrule\end{tabular}\end{center}
Circuit service envelopes meet the A1 limit of 100 m$^2$ and A2 limit of 50 m$^2$. Envelopes may overlap; they are not summed as office area. Connection and containment details follow.
\clearpage
% !TeX root = main.tex
\subsection*{FCU and containment details}
Each FB01--FB09 floorbox and SO01--SO03 twin socket shares one upstream 13 A FCU. FS01 uses a 10 A DP FCU and accessible earthed socket for WH01; retain the appliance cord and plug. Use accessible terminals rated for the 6/10 mm$^2$ main conductors. Place each FCU before its 4 mm$^2$, 0.5 m local spur; switch N with L and keep PE continuous. The demand basis is 13 A per socket assembly and 6.364 A for WH01.

Keep separate Method B containment banks: LC-A--G (seven circuits, $C_g=0.54$); P1/P2 and P3/P4 (four ring legs per bank, 0.65); P5/P6/P8 (three radials, 0.70); P7 alone (1.00); incoming feeder alone (1.00). Overlapping plan routes indicate corridors, not shared enclosures.

Under EMSD A5(3), Note 2, provide horizontal cable clearance between banks exceeding twice the largest adjacent cable diameter. Space enclosures to maintain that clearance; partitions alone do not provide thermal separation. Keep each circuit's L/N/CPC together and clear of thermal insulation. Recalculate grouping factors if this arrangement changes.

\subsection*{Protection and board arrangement}
The DB schematic shows NSX100H TM63D four-pole protection at MS-O1 before the CLP meter, 4-core 35 mm$^2$ plus CPC16, a 100 A four-pole DB isolator, independent L1/L2/L3/N/PE bars and 15 separate 1P+N Type A 30 mA RCBOs. Each circuit has a dedicated RCBO neutral and labelled phase; PE is unswitched.

Check driver inrush and protective-device coordination against the Task 3 requirements.
\subsection*{Emergency lighting assumption}
Landlord emergency lighting and exit signage use independent essential supplies, separated from normal lighting under EMSD 6B(4). These feeds bypass tenant normal-lighting switches and RCBO load sides. Emergency-system loads are therefore outside this office normal-supply DB.
